\section{Introduction}

Quadratic equations occupy an unusual position in mathematics: they belong to elementary algebra, yet they already contain several ideas that reappear in more advanced subjects. Their solution requires algebraic manipulation, the distinction between reversible and nonreversible transformations, the structure of number systems, and the interpretation of polynomial zeros. Their graphs lead directly to the geometry of the parabola, conic sections, symmetry, extrema, and tangency. Their negative discriminant forces an enlargement of the real numbers to the complex numbers, while their factorization anticipates the Fundamental Theorem of Algebra. In calculus, the same polynomial becomes a first model for derivatives, tangent lines, stationary points, and optimization. For this reason, a serious treatment of quadratic equations is not merely a rehearsal of a formula; it is a compact introduction to the way algebra, geometry, analysis, and mathematical modeling interact.

The historical development of the subject is correspondingly broad. Old Babylonian mathematical texts contain procedures for problems that, in modern notation, reduce to quadratic equations, although these procedures were expressed rhetorically and were embedded in numerical and geometric problems rather than in symbolic algebra \cite{robson2008,katz2009}. Greek mathematics treated geometrically related problems through constructions and the theory later described as ``application of areas.'' In India, Brahmagupta developed explicit rules for quadratic equations in the seventh century, and later authors, including \=Sr\={i}dhara and Bh\={a}skara II, refined algebraic techniques and discussed positive, negative, and irrational quantities with a sophistication that is essential to the history of the subject \cite{plofker2009,colebrooke1817}. In the Islamic mathematical tradition, al-Khw\={a}rizm\={i}'s ninth-century treatise systematically classified linear and quadratic equations and justified their solution through geometric arguments equivalent to completing the square \cite{rosen1831,swetzkatz2011}. Renaissance and early-modern symbolic algebra, particularly the work of Vi\`ete, Girard, and Descartes, gradually produced the notation in which the modern theory is now expressed \cite{boyer2011,katz2009}.

A terminology point is worth making at the outset. In Brazilian school mathematics the quadratic formula is often called the ``Bhaskara formula.'' This name is not standard in the international literature, where \emph{quadratic formula} is the usual term, and the modern symbolic formula should not be attributed to a single author. The mathematical ideas that led to it developed across several civilizations over many centuries. Bh\={a}skara II is an important figure in this history, but Brahmagupta, al-Khw\={a}rizm\={i}, earlier Mesopotamian mathematics, and the later development of symbolic algebra are also indispensable parts of the story \cite{plofker2009,katz2009,boyer2011}.

For readers who want a strong elementary foundation, Lang's \emph{Basic Mathematics} and Gelfand and Shen's \emph{Algebra} are particularly useful because they emphasize structure rather than mechanical substitution \cite{lang1988,gelfandshen1993}. Coxeter's \emph{Introduction to Geometry} is a classical reference for geometric thinking and conics \cite{coxeter1969}. Apostol and Spivak provide rigorous transitions from elementary algebra to calculus \cite{apostol1967,spivak2008}. For historical context, Katz, Boyer and Merzbach, Robson, and Plofker give complementary accounts of the Mesopotamian, Greek, Indian, Islamic, and European traditions \cite{katz2009,boyer2011,robson2008,plofker2009}. The aim of this note is to combine these perspectives in a self-contained exposition that begins with elementary prerequisites and ends with calculus and physical applications.

%%%%%%%%%%%%%%%%%%%%%%%%%%%%%%%%%%%%%%%%%%%%%%%%%%%%%%%%%%%%%%%%%%%%%%%
